\documentclass[11pt]{article}
\usepackage[margin=1in]{geometry}
\usepackage[T1]{fontenc}
\usepackage[utf8]{inputenc}
\usepackage{lmodern}
\usepackage{setspace}
\usepackage{parskip}
\usepackage{hyperref}

\setstretch{1.05}

\title{CS6423 Project Checkpoint Report}
\author{TokenSmith Team}
\date{\today}

\begin{document}
\maketitle

\section*{Project Overview}
Our project goal is to upgrade TokenSmith into a stronger retrieval-augmented generation system for database textbook question answering. The target architecture described in the proposal combines multiple retrieval signals rather than relying on a single chunk lookup method. In particular, we are building a hybrid pipeline that uses structure-aware chunking, semantic retrieval with FAISS, lexical retrieval with BM25, textbook-index keyword retrieval, and a LightRAG-style knowledge graph retriever. These retrieval sources are fused by an ensemble ranker and then passed into the existing answer generation stage. The broader objective is to improve retrieval quality, especially for textbook concepts that are distributed across sections, definitions, and related terms, while still keeping the system inspectable, deterministic, and practical to evaluate with benchmark metrics such as Recall@k and NDCG. The project repository is available at \url{https://github.com/XanderYoon/TokenSmith}, and my contributions for each of the three major upgrade areas were developed on their own Git branches: \texttt{chunking}, \texttt{ranking}, and \texttt{graph}, before being integrated into the combined branch.

\section*{Current Progress}
The first major area of progress is chunking. The original pipeline used a simpler recursive splitter that did not reliably preserve subsection boundaries or paragraph structure. We updated extraction and preprocessing so that top-level sections remain intact while subsection headings, paragraph breaks, list structure, and page markers are preserved. On top of that, we implemented a structure-aware chunking strategy that emits subsection-aligned or paragraph-aligned chunks and only falls back to local splitting when a unit is too large. This work is now wired through the configuration layer and indexing pipeline, and the index builder persists richer chunk metadata including section information, heading context, page numbers, unit type, and chunk identifiers. As a result, chunk boundaries are more faithful to the structure of the source textbook, which gives the later retrieval stages higher-quality units to work with.

The second major area of progress is hybrid retrieval and ranking. We cleaned up retrieval-source configuration so that enabled retrievers are explicit rather than being implied indirectly by old defaults. The ranking configuration now supports active retrievers and normalized source weights in a clearer way, and the ensemble ranker operates only over sources that are actually enabled. The query path has been updated so that FAISS, BM25, and textbook-index retrieval each produce their own score dictionary before fusion. This keeps the retrieval signals separate and makes the ranking behavior easier to reason about and test. We also added stronger automated coverage around hybrid retrieval, API behavior, and end-to-end execution so the retrieval stack can be validated deterministically. At this stage, the retrieval flow is no longer just a collection of independent components; it behaves like a single configurable pipeline in which each enabled retriever contributes a controlled signal to final ranking.

The third major area of progress is the knowledge graph layer. We added graph-specific modules for extraction, persistence, and retrieval. During indexing, the system now performs deterministic heuristic extraction over chunks to identify entities and supported relations, aggregates them into a versioned and human-reviewable JSON graph artifact, and stores that artifact alongside the FAISS, BM25, chunk, source, and metadata files. At query time, the graph artifact is loaded with the other retrieval artifacts, and a graph retriever matches normalized query terms against graph nodes, expands through one-hop edges, and maps evidence back to chunk scores. Importantly, graph retrieval has been integrated as a fourth first-class retrieval source rather than as a separate ranking path, so its scores are fused through the same ensemble mechanism as the other retrievers. This means the current branch now reflects the intended retrieval architecture from the proposal rather than only a partial implementation. Together, these three strands of work represent meaningful progress toward the full CS6423 design, because they improve the quality of the retrieval units, strengthen score fusion, and add a structured representation of textbook concepts and relations.

\section*{Challenges and Observations}
Several implementation challenges emerged while connecting these pieces. One recurring issue was that apparently small preprocessing choices had large downstream effects. For example, the original extraction logic split too aggressively on numbered headings and flattened structure in a way that made subsection-aware chunking impossible. Once we improved the structure preserved during extraction, the chunking stage became much more meaningful. Another observation is that hybrid retrieval only works well when each source has a clearly defined contract. Keeping per-source scores separate before fusion turned out to be essential for both correctness and debuggability. This was especially important once graph retrieval was added, because it needed to behave like the other retrievers instead of introducing ad hoc ranking logic.

The graph layer introduced its own set of challenges. Before retrieval could use graph information, we needed a stable artifact format that was persistent, inspectable, and reloadable. That forced us to define node, edge, and chunk-link representations early. We also chose a deterministic heuristic extraction approach for this checkpoint because it is easier to test and does not depend on external services, but that choice brings limitations: entity normalization is still fairly lightweight, relation coverage is limited, and one-hop neighborhood retrieval may miss some longer-range connections that could matter for textbook queries. Another challenge has been meshing these retrieval-side improvements with the front-end GitHub work, since integration requires the interface and backend behavior to remain aligned while the underlying retrieval stack is changing across multiple branches. In addition, some of the broader benchmark and metric tooling has environment-sensitive dependencies such as remote models or API-backed graders, which means practical offline validation requires graceful fallbacks. Overall, the system is now much closer to the proposal architecture, but retrieval quality will still depend on careful tuning of chunking parameters, source weights, and graph extraction quality.

\section*{Next Steps}
In the next few weeks, we plan to focus on evaluation, tuning, and refinement rather than large architectural rewrites. First, we want to run more robust benchmarking of each of the three major upgrade areas independently so we can identify their respective impacts on RAG performance. That means measuring the retrieval impact of structure-aware chunking against the earlier recursive baseline, isolating the contribution of the hybrid ranking changes to Recall@k and NDCG, and evaluating how much the graph retriever improves final retrieval quality when added as a fourth source. Second, we plan to tune the retrieval weights and graph scoring coefficients using the benchmark harness so that graph retrieval contributes useful evidence without overwhelming the other sources. Third, we want to improve graph extraction and normalization so that more textbook concepts, aliases, and relations are captured reliably. That includes expanding the heuristic patterns, improving singular-plural and paraphrase handling, and adding better introspection into which nodes and edges contributed to a retrieved chunk.

We also plan to strengthen usability and reporting around the retrieval stack. That includes adding retrieval-debug instrumentation, making artifact validation and failure modes more explicit, and supporting easier manual inspection of graph contents after indexing. If time permits, we would also like to explore modest graph expansion beyond one-hop retrieval when benchmarks show that it improves retrieval without adding too much noise. The main goal for the next checkpoint is therefore not to redesign the system, but to tighten the current implementation into a well-measured graph-enhanced hybrid retriever whose performance and tradeoffs are clearly documented.

\end{document}
